% ============================================================
\setpartpreamble{\partintro
  {Eigene SQL-Ausdrücke, Aggregate, Laufzeit-Metadaten und Scaffolding: die Werkzeuge für
  sehr komplexe Plugins.}
  {
    \item Kapitel~\ref{ch:custom-sql-expression-nodes}: Eigene SQL-Ausdrucksknoten
    \item Kapitel~\ref{ch:aggregate-translators}: Aggregat-Translators
    \item Kapitel~\ref{ch:property-level-fluent-api-and}: Fluent API auf Property-Ebene und
    Laufzeit-Annotations
    \item Kapitel~\ref{ch:scaffolding-in-practice}: Scaffolding in der Praxis
  }}
\part{Tiefe Erweiterungen}\label{part:deep-extensions}

% ------------------------------------------------------------
\chapter{Eigene SQL-Ausdrucksknoten}\label{ch:custom-sql-expression-nodes}
\versionrisk{volatile}

Kapitel~\ref{ch:translating-methods-and-members} baut SQL aus vorhandenen Bausteinen:
Funktionen, Konstanten, Casts. Das reicht für die meisten Plugins. Ein \textbf{eigener
Ausdrucksknoten} (eine Ableitung von \code{Sql\allowbreak{}Expression}) wird nötig, wenn
mindestens eins davon zutrifft:

\begin{datatable}{H{0.80}H{1.20}}
  \textbf{Grund} & \textbf{Beispiel} \\
  \tablehead
  Die Syntax ist kein Funktionsaufruf & Operatoren wie \code{@@}, \code{<->},
  \code{IS DISTINCT FROM}, Schlüsselwörter wie \code{AT TIME ZONE} \\
  \tablerule
  Die Nullability-Semantik weicht ab & \code{IS DISTINCT FROM} ist nie NULL, auch wenn ein
  Operand NULL ist \\
  \tablerule
  Andere Teile sollen den Knoten wiedererkennen & Optimierungen, die nur für diesen Ausdruck
  gelten \\
\end{datatable}

Das Beispiel implementiert den null-sicheren Vergleich
\code{EF.\allowbreak{}Functions.\allowbreak{}Is\allowbreak{}Not\allowbreak{}Distinct\allowbreak{}From(\allowbreak{}a,\allowbreak{}
b)} und \code{Is\allowbreak{}Distinct\allowbreak{}From(\allowbreak{}a,\allowbreak{} b)}
in einem eigenen kleinen Paket, \code{Acme.\allowbreak{}Entity\allowbreak{}Framework\allowbreak{}Core.\allowbreak{}Functions.\allowbreak{}Postgre\allowbreak{}SQL}, aktiviert mit \code{Use\allowbreak{}Npgsql(\allowbreak{}cs,\allowbreak{} o => o.\allowbreak{}Use\allowbreak{}Acme\allowbreak{}Functions(\allowbreak{}))}.
Npgsql bietet den Vergleich nicht als Funktion an; er erfüllt alle drei Gründe.

\begin{note}[Ein Lehrbeispiel, bewusst ohne Mandantenbezug]
  Der Vergleich hat nichts mit Mandantenfähigkeit zu tun. Er liegt in einem eigenen
  Beispielpaket, weil ein eigener Knoten ein Plugin zwingt, zwei Einzeldienste des Providers zu
  ersetzen: den SQL-Generator und den Nullability-Prozessor. Bis Ausgabe~1.1 tat das
  Mandanten-Plugin das allein für dieses Beispiel und ließ sich dadurch schwerer mit anderen
  Plugins kombinieren. Ein Mandanten-Plugin braucht keinen eigenen Knoten; Funktionskomposition
  (Kapitel~\ref{ch:translating-methods-and-members}) reicht. Die Tests betreiben beide Plugins
  nebeneinander, und das pgvector-Beispiel am Ende dieses Kapitels ergänzt ein drittes, das vom
  Generator des Funktionen-Plugins ableitet.
\end{note}

\section{Anatomie eines Knotens}

\begin{datatable}{H{0.78}H{1.26}H{0.96}}
  \textbf{Mitglied} & \textbf{Zweck} & \textbf{Folge, wenn es fehlt} \\
  \tablehead
  \code{Visit\allowbreak{}Children} / \code{Update} & Visitors können Kinder ersetzen;
  unveränderte Kinder ergeben dieselbe Instanz & EF wirft beim ersten Visitor \\
  \tablerule
  \code{Equals} / \code{Get\allowbreak{}Hash\allowbreak{}Code} & Strukturvergleich für den
  Query-Cache und Optimierungen & Cache-Fehltreffer, falsche Deduplizierung \\
  \tablerule
  \code{Print} & Lesbare Darstellung in DebugView und Fehlermeldungen & Unklare Fehler \\
  \tablerule
  \code{Quote} & C\#-Code für vorkompilierte Abfragen und NativeAOT & Kein AOT-Support \\
  \tablerule
  SQL-Generator & Schreibt den Knoten als SQL & \textbf{Stillschweigend} fehlendes SQL \\
  \tablerule
  Nullability-Prozessor & Kennt die NULL-Semantik & Exception beim Übersetzen \\
\end{datatable}

Die beiden letzten Zeilen betreffen auch die Kombination mit anderen Plugins: Ersetzt ein
zweites Plugin den Generator oder den Prozessor, fehlt das SQL des Knotens still oder die
Abfrage wirft. \code{Composition\allowbreak{}Matrix\allowbreak{}Tests} zeigt beides
(Kapitel~\ref{ch:registering-services}).

% source: src/Acme.EntityFrameworkCore.Functions.PostgreSQL/Query/Expressions/DistinctFromExpression.cs
\begin{sourcecode}[Query/\allowbreak{}Expressions/\allowbreak{}Distinct\allowbreak{}From\allowbreak{}Expression.\allowbreak{}cs (\allowbreak{}gekürzt)]{csharp}
public class DistinctFromExpression : SqlExpression
{
    private static ConstructorInfo? _quotingConstructor;

    public DistinctFromExpression(SqlExpression left, SqlExpression right, bool negated,
        RelationalTypeMapping? typeMapping)
        : base(typeof(bool), typeMapping)
    {
        // ... ThrowIfNull ...
        Left = left;
        Right = right;
        IsNegated = negated;
    }

    public virtual SqlExpression Left { get; }
    public virtual SqlExpression Right { get; }
    public virtual bool IsNegated { get; }

    protected override Expression VisitChildren(ExpressionVisitor visitor)
        => Update((SqlExpression)visitor.Visit(Left), (SqlExpression)visitor.Visit(Right));

    public virtual DistinctFromExpression Update(SqlExpression left, SqlExpression right)
        => left == Left && right == Right
            ? this
            : new DistinctFromExpression(left, right, IsNegated, TypeMapping);

    public override Expression Quote()
        => New(
            _quotingConstructor ??= typeof(DistinctFromExpression).GetConstructor(
                [typeof(SqlExpression), typeof(SqlExpression), typeof(bool),
                 typeof(RelationalTypeMapping)])!,
            Left.Quote(), Right.Quote(), Constant(IsNegated),
            RelationalExpressionQuotingUtilities.QuoteTypeMapping(TypeMapping));

    protected override void Print(ExpressionPrinter expressionPrinter)
    {
        expressionPrinter.Visit(Left);
        expressionPrinter.Append(IsNegated ? " IS DISTINCT FROM " : " IS NOT DISTINCT FROM ");
        expressionPrinter.Visit(Right);
    }

    public override bool Equals(object? obj)
        => obj is DistinctFromExpression other
           && (ReferenceEquals(this, other)
               || base.Equals(other) && Left.Equals(other.Left)
                  && Right.Equals(other.Right) && IsNegated == other.IsNegated);

    public override int GetHashCode()
        => HashCode.Combine(base.GetHashCode(), Left, Right, IsNegated);
}
\end{sourcecode}

Der Translator erzeugt den Knoten und gleicht dabei die Type Mappings beider Seiten an: Hat
eine Seite ein Mapping (etwa eine Spalte), bekommt die andere dasselbe. So wird aus einer
Konstante das passende Literal:

% source: src/Acme.EntityFrameworkCore.Functions.PostgreSQL/Query/DistinctFromTranslator.cs
\begin{sourcecode}[Query/\allowbreak{}Distinct\allowbreak{}From\allowbreak{}Translator.\allowbreak{}cs (\allowbreak{}Auszug)]{csharp}
var definition = method.GetGenericMethodDefinition();
if (definition != IsNotDistinctFromMethod && definition != IsDistinctFromMethod)
    return null;

var left = arguments[1];                        // arguments[0] ist DbFunctions
var right = arguments[2];
var typeMapping = ExpressionExtensions.InferTypeMapping(left, right)
                  ?? _typeMappingSource.FindMapping(left.Type);

return new DistinctFromExpression(
    _sql.ApplyTypeMapping(left, typeMapping)!,
    _sql.ApplyTypeMapping(right, typeMapping)!,
    negated: definition == IsDistinctFromMethod,
    _typeMappingSource.FindMapping(typeof(bool)));
\end{sourcecode}

\section{Den SQL-Generator erweitern}

Der \code{Query\allowbreak{}Sql\allowbreak{}Generator} schreibt SQL per Visitor. Unbekannte
Knoten gibt er an
\code{Expression\allowbreak{}Visitor.\allowbreak{}Visit\allowbreak{}Extension} weiter, das
nur die Kinder besucht. Ohne eigenen Override entsteht also kein Fehler, sondern \textbf{SQL
ohne den Operator}. Das Funktionen-Plugin erbt vom Npgsql-Generator:

% source: src/Acme.EntityFrameworkCore.Functions.PostgreSQL/Query/Internal/FunctionsNpgsqlQuerySqlGenerator.cs
\begin{sourcecode}[Query/\allowbreak{}Internal/\allowbreak{}Functions\allowbreak{}Npgsql\allowbreak{}Query\allowbreak{}Sql\allowbreak{}Generator.\allowbreak{}cs (\allowbreak{}gekürzt)]{csharp}
public class FunctionsNpgsqlQuerySqlGenerator : NpgsqlQuerySqlGenerator
{
    public FunctionsNpgsqlQuerySqlGenerator(QuerySqlGeneratorDependencies dependencies,
        IRelationalTypeMappingSource typeMappingSource,
        bool reverseNullOrderingEnabled, Version postgresVersion)
        : base(dependencies, typeMappingSource, reverseNullOrderingEnabled, postgresVersion)
    {
    }

    protected override Expression VisitExtension(Expression extensionExpression)
        => extensionExpression switch
        {
            DistinctFromExpression e => VisitDistinctFrom(e),
            _ => base.VisitExtension(extensionExpression),
        };

    protected virtual Expression VisitDistinctFrom(DistinctFromExpression expression)
    {
        VisitOperand(expression.Left);
        Sql.Append(expression.IsNegated ? " IS DISTINCT FROM " : " IS NOT DISTINCT FROM ");
        VisitOperand(expression.Right);
        return expression;
    }

    // Der Knoten als Operand anderer Operatoren: immer in Klammern
    protected override bool RequiresParentheses(SqlExpression outerExpression,
        SqlExpression innerExpression)
        => innerExpression is DistinctFromExpression
           || base.RequiresParentheses(outerExpression, innerExpression);
}
\end{sourcecode}

Klammern braucht es in zwei Richtungen. \code{Visit\allowbreak{}Operand} klammert
zusammengesetzte Operanden \emph{im} Knoten. Für den Knoten \emph{als} Operand fragt EF
\code{Requires\allowbreak{}Parentheses}, und das kennt nur seine eigenen Knoten. Ohne den
Override schrieb das Plugin
\code{a IS DISTINCT FROM 'x' <> a IS DISTINCT FROM 'y'}; PostgreSQL bindet \code{<>} stärker
als \code{IS} und meldet einen Syntaxfehler. \code{Deep\allowbreak{}Query\allowbreak{}Tests}
prüft das SQL, \code{Advanced\allowbreak{}Tests} die Abfrage gegen PostgreSQL.

Der Generator wird über
\code{IQuery\allowbreak{}Sql\allowbreak{}Generator\allowbreak{}Factory} erzeugt, einen
Singleton-Einzeldienst. Das Plugin ersetzt ihn durch eine eigene Factory mit denselben
Abhängigkeiten wie die von Npgsql
(\code{Query\allowbreak{}Sql\allowbreak{}Generator\allowbreak{}Dependencies},
Type-Mapping-Source, \code{INpgsql\allowbreak{}Singleton\allowbreak{}Options}).

\section{Den Nullability-Prozessor erweitern}

Nach der Übersetzung optimiert EF NULL-Vergleiche. Dafür muss es von jedem Knoten wissen, ob
er NULL werden kann. Unbekannte Knoten landen in
\code{Visit\allowbreak{}Custom\allowbreak{}Sql\allowbreak{}Expression}, dessen
Basisimplementierung eine Exception wirft. Der Prozessor hängt an einer Kette aus drei
Klassen, die alle ersetzt werden müssen:

% source: src/Acme.EntityFrameworkCore.Functions.PostgreSQL/Query/Internal/FunctionsNpgsqlSqlNullabilityProcessor.cs
\begin{sourcecode}[Query/\allowbreak{}Internal/\allowbreak{}Functions\allowbreak{}Npgsql\allowbreak{}Sql\allowbreak{}Nullability\allowbreak{}Processor.\allowbreak{}cs (\allowbreak{}gekürzt)]{csharp}
public class FunctionsNpgsqlSqlNullabilityProcessor : NpgsqlSqlNullabilityProcessor
{
    // Konstruktor: (…ProcessorDependencies, …ProcessorParameters) : base(...)

    protected override SqlExpression VisitCustomSqlExpression(
        SqlExpression sqlExpression, bool allowOptimizedExpansion, out bool nullable)
    {
        if (sqlExpression is DistinctFromExpression distinctFrom)
        {
            // Kinder besuchen (sie können optimiert werden), ihre Nullability ist egal
            var left = Visit(distinctFrom.Left, out _)!;
            var right = Visit(distinctFrom.Right, out _)!;

            nullable = false;                   // IS [NOT] DISTINCT FROM ist nie NULL
            return distinctFrom.Update(left, right);
        }

        return base.VisitCustomSqlExpression(
            sqlExpression, allowOptimizedExpansion, out nullable);
    }
}

public class FunctionsNpgsqlParameterBasedSqlProcessor : NpgsqlParameterBasedSqlProcessor
{
    protected override Expression ProcessSqlNullability(
        Expression selectExpression, ParametersCacheDecorator parametersDecorator)
        => new FunctionsNpgsqlSqlNullabilityProcessor(Dependencies, Parameters)
            .Process(selectExpression, parametersDecorator);
}

// + FunctionsNpgsqlParameterBasedSqlProcessorFactory
//   : IRelationalParameterBasedSqlProcessorFactory
\end{sourcecode}

% source: src/Acme.EntityFrameworkCore.Functions.PostgreSQL/Infrastructure/FunctionsOptionsExtension.cs
\begin{sourcecode}[Infrastructure/\allowbreak{}Functions\allowbreak{}Options\allowbreak{}Extension.\allowbreak{}cs › Apply\allowbreak{}Services (\allowbreak{}Ergänzung)]{csharp}
// 4) Tiefe Erweiterung: eigener SQL-Ausdrucksknoten
services.Replace(ServiceDescriptor
    .Singleton<IQuerySqlGeneratorFactory, FunctionsNpgsqlQuerySqlGeneratorFactory>());
services.Replace(ServiceDescriptor
    .Scoped<IRelationalParameterBasedSqlProcessorFactory,
        FunctionsNpgsqlParameterBasedSqlProcessorFactory>());
\end{sourcecode}

\begin{important}[Der Preis eigener Knoten]
  Eigene Knoten verlangen das Ersetzen zweier Einzeldienste und das Erben von internen
  Npgsql-Klassen. Zwei Plugins, die beide eigene Knoten einführen, lassen sich nur
  kombinieren, wenn eines vom anderen erbt. Deshalb gilt: Funktionskomposition zuerst;
  eigene Knoten nur, wenn einer der drei Gründe vorliegt. Und die Ersetzung gehört an
  prominente Stelle in die Dokumentation.
\end{important}

\begin{note}[Experimentelle API]
  \code{Quote} und
  \code{Relational\allowbreak{}Expression\allowbreak{}Quoting\allowbreak{}Utilities} gehören
  zur experimentellen API für vorkompilierte Abfragen. Ihre Verwendung meldet der Compiler
  als EF9100; das Beispiel unterdrückt die Meldung projektweit und dokumentiert das in
  \code{Directory.\allowbreak{}Build.\allowbreak{}props}.
\end{note}

% source: tests/Acme.EntityFrameworkCore.Tenancy.Tests/DeepQueryTests.cs
\begin{sourcecode}[tests/\allowbreak{}.\allowbreak{}.\allowbreak{}.\allowbreak{}/\allowbreak{}Deep\allowbreak{}Query\allowbreak{}Tests.\allowbreak{}cs (\allowbreak{}Auszug)]{csharp}
[Fact]
public void IsDistinctFrom_in_projection_passes_nullability_processor()
{
    // Ohne eigenen Nullability-Prozessor würde EF hier eine Exception werfen
    using var ctx = Contexts.Create();

    var sql = ctx.Customers
        .Select(c => EF.Functions.IsDistinctFrom(c.Name, (string?)null))
        .ToQueryString();

    Assert.Contains("IS DISTINCT FROM", sql);
}
\end{sourcecode}

\section{Praxisbeispiel: Operatoren für Vektorsuche und pgvector}

Embeddings und Ähnlichkeitssuche bringen Operatoren in SQL, die keine Funktionen sind: pgvector
schreibt die L2-Distanz als \code{<->}, die Kosinus-Distanz als
\code{<=>} und das negative innere Produkt als
\code{<\#>}; die Volltextsuche verwendet \code{@@}.
\code{Sql\allowbreak{}Binary\allowbreak{}Expression} hilft hier nicht, denn ihr Operator ist
ein \code{Expression\allowbreak{}Type}, und für \code{<=>} gibt es
keinen. \code{Custom\allowbreak{}Operator\allowbreak{}Tests} zeigt deshalb
\code{EF.\allowbreak{}Functions.\allowbreak{}Cosine\allowbreak{}Distance(\allowbreak{}a,\allowbreak{}
b)} als eigenen Knoten \code{Vector\allowbreak{}Distance\allowbreak{}Expression} mit dem
Operator als Zeichenkette, mit Translator, Generator und Nullability-Prozessor. Es ist
Testcode: Eine Anwendung nimmt \code{Pgvector.\allowbreak{}Entity\allowbreak{}Framework\allowbreak{}Core},
das den Typ \code{vector}, sein Type Mapping und diese Übersetzungen mitbringt. Npgsql selbst
hat für solche Fälle \code{Pg\allowbreak{}Unknown\allowbreak{}Binary\allowbreak{}Expression},
einen Knoten mit beliebigem Operator und propagierender Nullability, der Generator und
Prozessor nicht ersetzen muss. Er liegt aber in einem \code{Internal}-Namespace (EF1001).

Drei Dinge entscheiden darüber, ob ein solcher Operator richtiges SQL ergibt.

\textbf{Type Mappings.} Der Translator gibt beiden Operanden dasselbe Mapping, abgeleitet von
der Seite, die eins hat, meist der Spalte. So geht der Parameter mit dem Speichertyp der Spalte
an die Datenbank:

% source: tests/Acme.EntityFrameworkCore.Tenancy.Tests/CustomOperatorTests.cs
\begin{sourcecode}[tests/\allowbreak{}.\allowbreak{}.\allowbreak{}.\allowbreak{}/\allowbreak{}Custom\allowbreak{}Operator\allowbreak{}Tests.\allowbreak{}cs ›\allowbreak{} Translator]{csharp}
if (method != CosineDistanceMethod)
    return null;

var left = arguments[1];                        // arguments[0] ist DbFunctions
var right = arguments[2];
var operandMapping = ExpressionExtensions.InferTypeMapping(left, right)
                     ?? _typeMappingSource.FindMapping(left.Type);

return new VectorDistanceExpression(
    _sql.ApplyTypeMapping(left, operandMapping)!,
    _sql.ApplyTypeMapping(right, operandMapping)!,
    "<=>",
    _typeMappingSource.FindMapping(typeof(double)));
\end{sourcecode}

Im Test ist die Spalte ein \code{float[]}, also \code{real[]}. pgvector definiert seine
Operatoren aber nur für \code{vector}; gegen eine echte Datenbank braucht es deshalb ein Type
Mapping mit diesem Speichertyp, wie es das Paket registriert. Das Ergebnis bekommt das Mapping
von \code{double}, also \code{double precision}.

\textbf{Rangfolge und Klammern.} EF kennt die Rangfolge eines fremden Operators nicht.
PostgreSQL ordnet selbst definierte Operatoren zwischen die arithmetischen und die
Vergleichsoperatoren ein: \code{a <=> b * 2} heißt
\code{a <=> (\allowbreak{}b * 2)}. Ohne Gegenmaßnahme schrieb der
Generator für \code{Cosine\allowbreak{}Distance(\allowbreak{}.\allowbreak{}.\allowbreak{}.\allowbreak{})
* 2 + 1} genau das. Die Lösung sind dieselben zwei Richtungen wie oben: Operanden im Knoten
klammern und den Knoten als Operand über
\code{Requires\allowbreak{}Parentheses} immer klammern. \code{Print} klammert genauso, damit
die DebugView zeigt, was die Datenbank liest.

% source: tests/Acme.EntityFrameworkCore.Tenancy.Tests/CustomOperatorTests.cs
\begin{sourcecode}[tests/\allowbreak{}.\allowbreak{}.\allowbreak{}.\allowbreak{}/\allowbreak{}Custom\allowbreak{}Operator\allowbreak{}Tests.\allowbreak{}cs ›\allowbreak{} Generator]{csharp}
public class VectorQuerySqlGenerator : FunctionsNpgsqlQuerySqlGenerator
{
    // ... Konstruktor, VisitExtension, VisitOperand

    protected virtual Expression VisitVectorDistance(VectorDistanceExpression expression)
    {
        VisitOperand(expression.Left);
        Sql.Append($" {expression.Operator} ");
        VisitOperand(expression.Right);
        return expression;
    }

    // Ohne Override: "a <=> b * 2.0", gelesen als a <=> (b * 2.0)
    protected override bool RequiresParentheses(SqlExpression outerExpression,
        SqlExpression innerExpression)
        => innerExpression is VectorDistanceExpression
           || base.RequiresParentheses(outerExpression, innerExpression);
}
\end{sourcecode}

Der Generator erbt vom Generator des Plugins, nicht von dem von Npgsql: Beide ersetzen
denselben Einzeldienst, und nur so bleiben beide Knoten übersetzbar. Das ist der Preis aus
dem Kasten oben, im Kleinen.

\textbf{Nullability.} Ein gewöhnlicher Operator ist NULL, sobald ein Operand NULL ist. Der
Prozessor meldet das weiter:

% source: tests/Acme.EntityFrameworkCore.Tenancy.Tests/CustomOperatorTests.cs
\begin{sourcecode}[tests/\allowbreak{}.\allowbreak{}.\allowbreak{}.\allowbreak{}/\allowbreak{}Custom\allowbreak{}Operator\allowbreak{}Tests.\allowbreak{}cs ›\allowbreak{} Nullability]{csharp}
if (sqlExpression is VectorDistanceExpression distance)
{
    // NULL hinein, NULL heraus
    var left = Visit(distance.Left, out var leftNullable)!;
    var right = Visit(distance.Right, out var rightNullable)!;

    nullable = leftNullable || rightNullable;
    return distance.Update(left, right);
}
\end{sourcecode}

Die Meldung hat sichtbare Folgen. Für
\code{!(\allowbreak{}Cosine\allowbreak{}Distance(\allowbreak{}.\allowbreak{}.\allowbreak{}.\allowbreak{})
< 0.5)} schreibt EF bei einer Pflichtspalte \code{>= 0.5}, bei einer optionalen Spalte ein
\code{CASE}, weil \code{NOT (\allowbreak{}NULL < 0.5)} in SQL NULL ist, in C\# aber
\code{true}. Ohne den Prozessor scheitert die Übersetzung. Die Tests prüfen das SQL aller
Fälle:

% source: tests/Acme.EntityFrameworkCore.Tenancy.Tests/CustomOperatorTests.cs
\begin{sourcecode}[tests/\allowbreak{}.\allowbreak{}.\allowbreak{}.\allowbreak{}/\allowbreak{}Custom\allowbreak{}Operator\allowbreak{}Tests.\allowbreak{}cs (\allowbreak{}Auszug)]{csharp}
var sql = ctx.Documents
    .Select(d => EF.Functions.CosineDistance(d.Embedding, query) * 2 + 1)
    .ToQueryString();

Assert.Contains("SELECT (d.\"Embedding\" <=> @query) * 2.0 + 1.0", sql);
\end{sourcecode}

% ------------------------------------------------------------
\chapter{Aggregat-Translators}\label{ch:aggregate-translators}
\versionrisk{medium}

Aggregate wie \code{Sum} oder \code{string\_\allowbreak{}agg} fassen mehrere Zeilen
zusammen, typischerweise nach \code{Group\allowbreak{}By}. Ein
\code{IAggregate\allowbreak{}Method\allowbreak{}Call\allowbreak{}Translator} bekommt statt
eines einzelnen Arguments eine \code{Enumerable\allowbreak{}Expression}, die beschreibt,
\emph{was} aggregiert wird:

\begin{datatable}{H{0.70}H{1.09}H{1.21}}
  \textbf{Eigenschaft} & \textbf{Bedeutung} & \textbf{Herkunft in LINQ} \\
  \tablehead
  \code{Selector} & Der aggregierte Ausdruck & \code{g.\allowbreak{}Select(\allowbreak{}o =>
  o.\allowbreak{}Total > 0)} \\
  \tablerule
  \code{Is\allowbreak{}Distinct} & Nur verschiedene Werte &
  \code{g.\allowbreak{}Select(\allowbreak{}.\allowbreak{}.\allowbreak{}.\allowbreak{}).\allowbreak{}Distinct(\allowbreak{})}
  \\
  \tablerule
  \code{Predicate} & Filter vor dem Aggregat & \code{g.\allowbreak{}Where(\allowbreak{}o =>
  o.\allowbreak{}Id > 10)} \\
  \tablerule
  \code{Orderings} & Reihenfolge (für reihenfolgeabhängige Aggregate) &
  \code{g.\allowbreak{}Order\allowbreak{}By(\allowbreak{}.\allowbreak{}.\allowbreak{}.\allowbreak{})}
  \\
\end{datatable}

Das Funktionen-Beispielpaket aus Kapitel~\ref{ch:custom-sql-expression-nodes} ergänzt
\code{EF.\allowbreak{}Functions.\allowbreak{}Every(\allowbreak{}values)} für PostgreSQLs
\code{bool\_\allowbreak{}and}; wie der Vergleich ist es ein Lehrbeispiel ohne
Mandantenbezug:

% source: src/Acme.EntityFrameworkCore.Functions.PostgreSQL/Query/EveryAggregateTranslator.cs
\begin{sourcecode}[Query/\allowbreak{}Every\allowbreak{}Aggregate\allowbreak{}Translator.\allowbreak{}cs (\allowbreak{}Auszug)]{csharp}
public virtual SqlExpression? Translate(MethodInfo method, EnumerableExpression source,
    IReadOnlyList<SqlExpression> arguments,
    IDiagnosticsLogger<DbLoggerCategory.Query> logger)
{
    if (method != EveryMethod || source.Selector is not SqlExpression selector)
        return null;

    // Npgsqls Factory setzt DISTINCT, FILTER (WHERE ...) und ORDER BY im Aggregat um
    if (_sql is not NpgsqlSqlExpressionFactory npgsql)
        return null;

    var boolean = _typeMappingSource.FindMapping(typeof(bool))!;

    // bool_and liefert für eine leere Menge NULL, Enumerable.All liefert true
    return _sql.Coalesce(
        npgsql.AggregateFunction("bool_and", [selector], source,
            nullable: true, argumentsPropagateNullability: [false],
            typeof(bool), boolean),
        _sql.Constant(true, boolean),
        boolean);
}
\end{sourcecode}

\begin{sourcecode}[Erzeugtes SQL für g.\allowbreak{}Where(\allowbreak{}o => o.\allowbreak{}Id > 10).\allowbreak{}Select(\allowbreak{}o => o.\allowbreak{}Total > 0)]{sql}
SELECT COALESCE(bool_and(o."Total" > 0.0) FILTER (WHERE o."Id" > 10), TRUE)
FROM "Orders" AS o
GROUP BY o."TenantId"
\end{sourcecode}

\begin{important}[Semantik vor Syntax]
  Ein Aggregat-Translator muss die Semantik der .NET-Seite treffen, nicht nur die Syntax der
  Datenbank. Typische Fallen: leere Mengen (NULL statt 0, true oder Leerstring), NULL-Werte
  in der Eingabe und Prädikate. Npgsql selbst gleicht etwa \code{string.\allowbreak{}Join}
  mit \code{COALESCE} an \code{string\_\allowbreak{}agg} an; das Beispiel folgt diesem
  Muster.
\end{important}

Datenbanken ohne \code{FILTER}-Klausel (etwa SQL Server) emulieren Prädikate mit \code{CASE
WHEN prädikat THEN wert END}. Ein provider-neutraler Aggregat-Translator muss beides kennen;
dieser nutzt die Npgsql-Factory und liegt deshalb im PostgreSQL-spezifischen Beispielpaket.

% ------------------------------------------------------------
\chapter{Fluent API auf Property-Ebene und Laufzeit-Annotations}\label{ch:property-level-fluent-api-and}
\versionrisk{volatile}

\section{Property-Builder}

Bisher setzt das Plugin voraus, dass die Mandanten-Property \code{Tenant\allowbreak{}Id}
heißt und \code{ITenant\allowbreak{}Owned} implementiert ist. Mit einer Fluent-API-Methode
auf \code{Property\allowbreak{}Builder} geht es ohne beides, etwa bei bestehenden Modellen
mit einer Property \code{Mandant\allowbreak{}Id} oder einer expliziten
Interface-Implementierung:

\begin{sourcecode}[Anwendung]{csharp}
modelBuilder.Entity<Contract>().Property(c => c.Mandant).IsTenantProperty();
\end{sourcecode}

% source: src/Acme.EntityFrameworkCore.Tenancy/Metadata/TenancyPropertyBuilderExtensions.cs
\begin{sourcecode}[Metadata/\allowbreak{}Tenancy\allowbreak{}Property\allowbreak{}Builder\allowbreak{}Extensions.\allowbreak{}cs (\allowbreak{}Kernpaket)]{csharp}
public static PropertyBuilder IsTenantProperty(
    this PropertyBuilder builder, bool isTenant = true)
{
    ArgumentNullException.ThrowIfNull(builder);

    if (isTenant)
        builder.Metadata.SetAnnotation(TenancyAnnotationNames.IsTenantProperty, true);
    else
        builder.Metadata.RemoveAnnotation(TenancyAnnotationNames.IsTenantProperty);

    return builder;
}

public static PropertyBuilder<TProperty> IsTenantProperty<TProperty>(
    this PropertyBuilder<TProperty> builder, bool isTenant = true)
    => (PropertyBuilder<TProperty>)IsTenantProperty((PropertyBuilder)builder, isTenant);

public static IConventionPropertyBuilder? IsTenantProperty(
    this IConventionPropertyBuilder builder,
    bool isTenant = true,
    bool fromDataAnnotation = false)
    => builder.HasAnnotation(TenancyAnnotationNames.IsTenantProperty,
           isTenant ? true : null, fromDataAnnotation) is null
        ? null
        : builder;
\end{sourcecode}

Die Methode liegt im Kernpaket, weil sie nur Metadaten setzt. Die Conventions aller
Provider-Pakete suchen die Mandanten-Property über dieselbe Kernfunktion:

\begin{sourcecode}[Metadata/\allowbreak{}Tenancy\allowbreak{}Entity\allowbreak{}Type\allowbreak{}Extensions.\allowbreak{}cs (\allowbreak{}Kernpaket,\allowbreak{} gekürzt, vereinfacht)]{csharp}
public static IReadOnlyProperty? FindTenantPropertyCandidate(
    this IReadOnlyEntityType entityType)
{
    var marked = entityType.GetProperties()
        .Where(p => p[TenancyAnnotationNames.IsTenantProperty] is true)
        .ToList();

    if (marked.Count > 1)
        throw new InvalidOperationException("... mehrere Mandanten-Properties ...");
    if (marked.Count == 1)
        return marked[0];

    return typeof(ITenantOwned).IsAssignableFrom(entityType.ClrType)
        ? entityType.FindProperty(nameof(ITenantOwned.TenantId))
        : null;
}
\end{sourcecode}

\begin{datatable}{H{0.92}H{1.10}H{0.98}}
  \textbf{Builder} & \textbf{Typische Plugin-Methoden} & \textbf{Hinweis} \\
  \tablehead
  \code{Property\allowbreak{}Builder} / \code{<\allowbreak{}T>} & Spaltenoptionen, Rollen
  von Properties & Hier: \code{Is\allowbreak{}Tenant\allowbreak{}Property} \\
  \tablerule
  \code{Index\allowbreak{}Builder} / \code{<\allowbreak{}T>} & Index-Methoden,
  Operator-Klassen, Filter & Annotation an \code{Metadata} (IMutableIndex) \\
  \tablerule
  \code{Key\allowbreak{}Builder} & Schlüsseloptionen (Clustering) & Wie oben \\
  \tablerule
  \code{Complex\allowbreak{}Type\allowbreak{}Property\allowbreak{}Builder} & Properties in
  Complex Types & Nur anbieten, wenn die Option dort Sinn ergibt
  (Kapitel~\ref{ch:complex-types-and-json}) \\
  \tablerule
  \code{Model\allowbreak{}Builder} & Modellweite Optionen & Annotation am Modell \\
\end{datatable}

\section{Laufzeit-Annotations}

Der SaveChanges-Interceptor muss zur Laufzeit wissen, welche Entitäten isoliert sind und wie
ihre Mandanten-Property heißt. Die Convention schreibt das in eine Annotation am Entity-Typ;
der Interceptor liest sie aus dem \emph{Laufzeitmodell}:

\begin{sourcecode}[Convention und Kernpaket (\allowbreak{}Auszug)]{csharp}
// Convention (Design-Time-Modell)
entityType.Builder.HasAnnotation(TenancyAnnotationNames.TenantProperty, property.Name);

// Interceptor (Laufzeitmodell)
public static string? GetTenantPropertyName(this IReadOnlyEntityType entityType)
    => entityType.GetRootType()[TenancyAnnotationNames.TenantProperty] as string;
\end{sourcecode}

Zwei Fragen entscheiden, ob das funktioniert. Landet die Annotation im Laufzeitmodell? Ja:
\code{Runtime\allowbreak{}Model\allowbreak{}Convention} entfernt nur EFs eigene
Design-Time-Annotations und übernimmt alle anderen. Landet sie im Compiled Model? Ja:
Einfache Werte schreibt EF automatisch
(Kapitel~\ref{ch:compiled-models-trimming-and-nativeaot}). Der eigene
\code{ICSharp\allowbreak{}Runtime\allowbreak{}Annotation\allowbreak{}Code\allowbreak{}Generator}
des PostgreSQL-Pakets hat deshalb die umgekehrte Aufgabe: Er entfernt die RLS-Annotations, die nur Migrations
brauchen:

\begin{sourcecode}[Design/\allowbreak{}Internal/\allowbreak{}Tenancy\allowbreak{}CSharp\allowbreak{}Runtime\allowbreak{}Annotation\allowbreak{}Code\allowbreak{}Generator.\allowbreak{}cs (\allowbreak{}gekürzt, vereinfacht)]{csharp}
public class TenancyCSharpRuntimeAnnotationCodeGenerator
    : NpgsqlCSharpRuntimeAnnotationCodeGenerator
{
    // Konstruktor: (CSharpRuntimeAnnotationCodeGeneratorDependencies,
    //               RelationalCSharpRuntimeAnnotationCodeGeneratorDependencies) : base(...)

    public override void Generate(IEntityType entityType,
        CSharpRuntimeAnnotationCodeGeneratorParameters parameters)
    {
        if (!parameters.IsRuntime)
        {
            foreach (var key in parameters.Annotations.Keys
                         .Where(k => k.StartsWith(RlsAnnotationNames.Prefix,
                                                  StringComparison.Ordinal))
                         .ToList())
            {
                parameters.Annotations.Remove(key);
            }
        }

        base.Generate(entityType, parameters);
    }
}
\end{sourcecode}

Wird das Compiled Model für NativeAOT erzeugt
(\code{For\allowbreak{}Native\allowbreak{}Aot}), schreibt EF zusätzlich das relationale
Modell. Der Annotation-Provider hat die RLS-Annotations auf die Tabellen kopiert, sie
tauchen dort also ein zweites Mal auf. Der Generator wendet deshalb denselben Filter in
einer Überschreibung von \code{Generate(\allowbreak{}ITable,\allowbreak{}
.\allowbreak{}.\allowbreak{}.\allowbreak{})} an.

\begin{datatable}{H{1.36}H{0.88}H{0.88}H{0.88}}
  \textbf{Annotation} & \textbf{Design-Time} & \textbf{Laufzeit} & \textbf{Compiled Model} \\
  \tablehead
  \code{Acme:Rls:*} (Policies) & ja (Migrations) & ja, ungenutzt & gefiltert \\
  \tablerule
  \code{Acme:Tenancy:Tenant\allowbreak{}Property} & ja & ja (Interceptor) & ja \\
  \tablerule
  \code{Acme:Tenancy:Is\allowbreak{}Tenant\allowbreak{}Property} & ja (Convention) & ja,
  ungenutzt & ja (klein) \\
\end{datatable}

\begin{note}[Laufzeit-Annotation und Runtime-Annotation]
  „Laufzeit-Annotation“ heißt in diesem Buch: eine gewöhnliche Annotation, die zur Laufzeit
  gelesen wird. EF kennt zusätzlich eine eigene Runtime-Annotation-API
  (\code{Add\allowbreak{}Runtime\allowbreak{}Annotation},
  \code{Get\allowbreak{}Or\allowbreak{}Add\allowbreak{}Runtime\allowbreak{}Annotation\allowbreak{}Value}
  auf \code{IAnnotatable}). Sie legt Werte erst am fertigen Modell ab, etwa abgeleitete
  Daten, die nicht bei jedem Zugriff neu berechnet werden sollen. Solche Werte erscheinen nie
  in Snapshot oder Migrations. Für Konfiguration, die ein Plugin beim Modellbau festlegt, ist
  die gewöhnliche Annotation der richtige Weg.
\end{note}

\begin{tip}
  Für jede Annotation sollte klar sein, wer sie wann liest. Nur zur Design-Zeit gelesene
  Annotations gehören nicht ins Compiled Model; zur Laufzeit gelesene müssen hinein und
  brauchen einen Test gegen das generierte Modell. Der Compiled-Model-Test des Beispiels
  prüft beides.
\end{tip}

% ------------------------------------------------------------
\chapter{Scaffolding in der Praxis}\label{ch:scaffolding-in-practice}
\versionrisk{volatile}

Kapitel~\ref{ch:scaffolding-reverse-engineering} nennt die beiden Design-Time-Dienste, über
die ein Plugin seine Schema-Objekte beim Scaffolding erkennt. Dieses Kapitel zeigt die
Umsetzung. Der Weg führt über drei Stationen, von denen das Plugin zwei bereits kennt:

\begin{flowdiagram}
  \flowstep{IDatabaseModelFactory}{Liest das Schema in ein DatabaseModel; das Plugin hängt
  RLS-Daten als Annotations an die Tabellen}
  \flowstep{Scaffolding-Model-Factory}{EF überträgt Tabellen-Annotations auf die
  Entity-Typen}
  \flowstep{IAnnotationCodeGenerator}{Schreibt sie als HasRowLevelSecurity/HasPolicy
  (Kapitel~\ref{ch:design-time-services})}
\end{flowdiagram}

% source: src/Acme.EntityFrameworkCore.Tenancy.PostgreSQL/Scaffolding/TenancyDatabaseModelFactory.cs
\begin{sourcecode}[Scaffolding/\allowbreak{}Tenancy\allowbreak{}Database\allowbreak{}Model\allowbreak{}Factory.\allowbreak{}cs (\allowbreak{}gekürzt)]{csharp}
public class TenancyDatabaseModelFactory : NpgsqlDatabaseModelFactory
{
    public TenancyDatabaseModelFactory(
        IDiagnosticsLogger<DbLoggerCategory.Scaffolding> logger)
        : base(logger)
    {
    }

    public override DatabaseModel Create(
        DbConnection dbConnection, DatabaseModelFactoryOptions options)
    {
        var model = base.Create(dbConnection, options);
        // ... Verbindung öffnen, falls nötig ...
        ReadRowLevelSecurity(dbConnection, model);    // pg_class.relrowsecurity / ...force...
        ReadPolicies(dbConnection, model);            // pg_policies
        return model;
    }

    protected virtual void ReadPolicies(DbConnection connection, DatabaseModel model)
    {
        // SELECT schemaname, tablename, policyname, permissive, roles::text[], cmd,
        //        qual, with_check FROM pg_catalog.pg_policies
        // ...
        var policy = new RlsPolicy(
            Name: reader.GetString(2),
            Command: Enum.Parse<PolicyCommand>(reader.GetString(5), ignoreCase: true),
            Permissive: reader.GetString(3) == "PERMISSIVE",
            Roles: roles is ["public"] ? null : roles,    // "public" = Default ohne TO
            Using: reader.IsDBNull(6) ? null : reader.GetString(6),
            WithCheck: reader.IsDBNull(7) ? null : reader.GetString(7));

        table[RlsAnnotationNames.PolicyPrefix + policy.Name] = policy.Serialize();
    }
}
\end{sourcecode}

Registriert wird die Factory wie der Code-Generator in den Design-Time-Services mit
\code{Remove\allowbreak{}All}, aber mit \code{Add\allowbreak{}Scoped}: EF und Npgsql
registrieren die Database-Model-Factory als Scoped, und sie hängt vom ebenfalls als Scoped registrierten
Scaffolding-Logger ab. Das Datenformat ist dasselbe wie bei der Fluent API; deshalb
funktionieren Snapshot, Migrations und Code-Generator ohne weitere Anpassung.

\begin{note}[Normalisierte Ausdrücke]
  PostgreSQL speichert Policy-Ausdrücke in normalisierter Form, etwa
  \code{(\allowbreak{}(\allowbreak{}\textquotedbl{}Tenant\allowbreak{}Id\textquotedbl{} =
  (\allowbreak{}NULLIF(\allowbreak{}current\_\allowbreak{}setting(\allowbreak{}'app.\allowbreak{}tenant\_\allowbreak{}id'::text,\allowbreak{}
  true),\allowbreak{} ''::text))::uuid))}. Gescaffoldete Policies enthalten genau diesen
  Text. Das ist korrekt, sieht aber anders aus als die von der Convention erzeugte Form.
  Teams, die nach dem Scaffolding auf Migrations umsteigen, sollten die Policies einmal
  vereinheitlichen, sonst erzeugt die erste Migration ein DROP/CREATE mit gleichem Inhalt.
\end{note}

% source: tests/Acme.EntityFrameworkCore.Tenancy.IntegrationTests/AdvancedTests.cs
\begin{sourcecode}[tests/\allowbreak{}.\allowbreak{}.\allowbreak{}.\allowbreak{}Integration\allowbreak{}Tests/\allowbreak{}Advanced\allowbreak{}Tests.\allowbreak{}cs (\allowbreak{}Auszug)]{csharp}
[Fact]
public void Scaffolding_reads_rls_and_policies()
{
    using var ctx = ContextFor(null);
    var factory = new TenancyDatabaseModelFactory(
        ctx.GetService<IDiagnosticsLogger<DbLoggerCategory.Scaffolding>>());

    var model = factory.Create(
        _fixture.OwnerConnectionString, new DatabaseModelFactoryOptions());
    var orders = model.Tables.Single(t => t.Name == "Orders");

    Assert.Equal(true, orders[RlsAnnotationNames.Enabled]);
    Assert.Equal(true, orders[RlsAnnotationNames.Forced]);
}
\end{sourcecode}
